\chapter{Value-Based RL: Q-Learning and DQN}
\section{Learning Goals and Roadmap}
Q-Learning $\rightarrow$ neural Q-functions $\rightarrow$ DQN
$\rightarrow$ Double DQN. We study action selection, data, targets,
losses, gradient flow, and limitations.

\section{Why Do We Need It?}
The objective is still to maximize expected discounted return.
A state value $V^\pi(s)$ evaluates a state, but alone does not rank
actions without additional knowledge of rewards and dynamics.
An action value evaluates a particular first action:
\begin{equation}
Q^\pi(s,a)=\E_\pi\left[\sum_{k=0}^{\infty}\gamma^k r_{t+k}
\mid s_t=s,a_t=a\right].
\end{equation}
After that first action, the agent follows $\pi$. The optimal
action value $Q^*(s,a)$ assumes optimal subsequent behavior.
For finite discrete actions, a greedy policy selects
$\arg\max_a Q(s,a)$. This is the best action according to the
current estimate, not necessarily the truly optimal action.

\section{Tabular Q-Learning}
Store one estimate for each state-action pair. For a sampled
transition $(s,a,r,s',d)$, with $d=1$ for genuine termination,
\begin{align}
y &= r+\gamma(1-d)\max_{a'}Q(s',a'),\\
\delta &= y-Q(s,a),\\
Q(s,a)&\leftarrow Q(s,a)+\alpha\delta.
\end{align}
The Bellman optimality equation motivates this sampled target:
\begin{equation}
Q^*(s,a)=\E_{r,s'\mid s,a}
[r+\gamma(1-d)\max_{a'}Q^*(s',a')].
\end{equation}
Example: $Q(s,a)=3$, $r=2$, $\gamma=0.9$, and next-state
values $(4,6)$. Then $y=7.4$, $\delta=4.4$, and with
$\alpha=0.1$ the updated entry is $3.44$.
The maximum expresses optimal subsequent behavior; averaging
action values would correspond to a different subsequent policy.

\subsection{Exploration and Off-Policy Learning}
An epsilon-greedy behavior policy chooses a random action with
probability $\epsilon$ and a greedy action otherwise.
The update still uses the greedy next-state value even if the
next action actually taken is exploratory. Thus data collection
and the target policy can differ: Q-learning is off-policy.
If next-state values are $(4,6)$ and exploration chooses the first
action, the previous transition's bootstrap value remains $6$.

\section{Models and Inputs/Outputs}
A neural Q-function replaces a table when states are high-dimensional
or too numerous. A common discrete-action architecture maps
$s\in\R^{n_s}$ to a vector in $\R^{n_a}$, with one Q-value per action.
These are return estimates, not probabilities; no softmax is needed.
For a batch, states have shape $[B,n_s]$ and outputs $[B,n_a]$.
For each transition, select the value of the action actually executed,
not the largest current-state output.
Example: outputs $(2,3)$ and stored action index $0$ give prediction $2$.
Shared parameters enable generalization but also make updates affect
other states and actions.

\section{DQN: Replay and Target Networks}
DQN combines a neural Q-function with experience replay and a slowly
updated target network.
\begin{description}[style=nextline]
\item[Online network $Q_\theta$] Selects interaction actions and produces
current predictions. Its parameters are updated by gradient descent.
\item[Target network $Q_{\bar\theta}$] Produces bootstrap estimates.
It is initialized from the online network and periodically refreshed
by $\bar\theta\leftarrow\theta$ in the hard-update version studied here.
\item[Replay buffer] Stores observed transitions $(s,a,r,s',d)$.
Random minibatches reduce temporal correlation and reuse experience.
Stored transitions remain fixed; targets are computed when sampled.
\end{description}
\begin{align}
y_i&=r_i+\gamma(1-d_i)\max_{a'}Q_{\bar\theta}(s'_i,a'),\\
L(\theta)&=\frac1B\sum_i
\left[Q_\theta(s_i,a_i)-\operatorname{stopgrad}(y_i)\right]^2.
\end{align}
Squared error illustrates the principle; Huber loss is also common.
Gradient flows through the current prediction only.
Detaching a target prevents differentiation through it in that update,
but does not make its numerical value permanently fixed.

\subsection{Complete Training Flow}
\begin{enumerate}
\item Use the online network with epsilon-greedy exploration to act.
\item Observe the transition and store it in replay.
\item Randomly sample a minibatch, possibly containing old experience.
\item Compute executed-action predictions with the online network.
\item Compute next-state targets with the target network, without gradients.
\item Backpropagate the loss and update the online network.
\item At the chosen interval, copy online parameters into the target network.
\end{enumerate}
Interaction and optimization need not use the same transition.
Updating Q-values implicitly changes the greedy policy: values $(5,3)$
choose left, while later values $(5,7)$ choose right, without an actor update.

\subsection{Minimal Code Data Flow}
\begin{verbatim}
# states, next_states: [B, obs_dim]
# actions: [B, 1], integer indices
# rewards, terminated: [B, 1]
q_pred = online_net(states).gather(1, actions)
with torch.no_grad():
    next_best = target_net(next_states).max(
        dim=1, keepdim=True).values
    target = rewards + gamma * (
        1 - terminated.float()) * next_best
loss = ((q_pred - target) ** 2).mean()
optimizer.zero_grad()
loss.backward()
optimizer.step()
# Periodically: target_net.load_state_dict(online_net.state_dict())
\end{verbatim}
This is schematic code. Handle environment resets and truncations
correctly: a time-limit truncation is not automatically a genuine
terminal state, and bootstrapping may require its final observation.

\begin{pitfall}
For ordinary DQN, a fixed transition has a fixed target between
hard target-network updates, assuming deterministic network outputs.
When the target network is refreshed, the target can change even
though the observed reward has not changed.
\end{pitfall}

\section{Double DQN: Separate Selection and Evaluation}
Taking a maximum over noisy estimates can cause overestimation.
If both true action values are $5$, estimates alternating between
$(4,6)$ and $(6,4)$ are individually unbiased, but their maximum
is always $6$.
Ordinary DQN uses the target network both to select the largest
next-state action value and to evaluate that action.
Double DQN instead uses:
\begin{align}
a_i^*&=\arg\max_{a'}Q_\theta(s'_i,a'),\\
y_i^{Double}&=r_i+\gamma(1-d_i)Q_{\bar\theta}(s'_i,a_i^*).
\end{align}
The existing two networks and their update schedule can remain the same;
the change is the construction of the target.
For online values $(8,3)$ and target values $(4,7)$, DQN bootstraps
with $7$, while Double DQN selects the first action online and evaluates
it as $4$ using the target network.
Separating selection and evaluation reduces max-related overestimation,
but does not guarantee removal of all bias; network errors remain correlated.
If both networks are identical, their targets coincide.

\begin{verbatim}
with torch.no_grad():
    next_actions = online_net(next_states).argmax(
        dim=1, keepdim=True)                 # [B, 1]
    next_q = target_net(next_states).gather(
        1, next_actions)                     # [B, 1]
    target = rewards + gamma * (
        1 - terminated.float()) * next_q
\end{verbatim}
The prediction and loss are unchanged.
Stored actions index actual past behavior at the current state.
Next actions are freshly selected to construct next-state targets.
Both operations above occur without gradients.

\begin{pitfall}
Double DQN targets can change even while target-network parameters
are fixed, because the online-selected action can change.
If target values remain $(4,7)$ while online values change from
$(8,3)$ to $(2,9)$, the bootstrap value changes from $4$ to $7$.
This selection concerns the training target. Environment interaction
still ordinarily uses online epsilon-greedy action selection.
\end{pitfall}

\section{Links to Other Methods and Limitations}
Yesterday's stochastic actor--state-value critic learns a policy
distribution and a scalar $V(s)$. The critic helps train the actor
through advantage estimates. DQN learns action values and derives
a greedy policy directly. Actor-critic critics may also be Q-functions,
and actors need not always be stochastic.
PPO reuses a recent rollout with fixed batch advantages and return
targets over optimization epochs; DQN reuses replay transitions
and recomputes bootstrap targets when sampled.

Standard DQN enumerates finite actions. Continuous actions would require
solving $\arg\max_a Q(s,a)$ over a continuous space at each decision.
The next chapter introduces actors that directly output continuous
actions and learn to increase their critic's estimated Q-value.
For a deterministic actor, this motivates
$J(\theta)=\E_s[Q_\phi(s,\mu_\theta(s))]$.

Neural approximation, bootstrapping, and off-policy learning together
form the deadly triad: their combination can cause instability or
divergence. Replay, target networks, and Double DQN improve practical
learning but do not provide general convergence guarantees for
nonlinear neural Q-learning. Exploration and data coverage also matter.

\section{Review Questions}
\begin{enumerate}
\item Why can a Q-function directly rank actions while a state value alone cannot?
\item Why does Q-learning take a next-action maximum instead of an average?
\item Why can an exploratory transition train a greedy target policy?
\item Which current-state output enters the loss, and why?
\item How does no-gradient target computation differ from freezing its values?
\item Why can the same replay transition receive a new DQN target?
\item What are the separate roles of replay and the target network?
\item In Double DQN, which network selects and which evaluates the next action?
\item Can a Double DQN target change before the target network is refreshed?
\item Why is continuous-action maximization harder, and how might an actor help?
\end{enumerate}

\section{Learning Status and Source Notes}
First study pass completed on October 13, 2026.
Clarified distinctions: executed action versus target action;
ordinary DQN's two networks versus Double DQN's selection/evaluation split;
fixed transitions versus recomputed targets; and implicit greedy-policy
improvement versus explicit actor learning.
Course basis: supplied Lectures 4 and 5 on value-based methods and
deep Q-functions, supplemented by our guided explanations and examples.
This chapter is a first-pass study note, not an exhaustive account of
all lecture topics or a verbatim slide transcription.
